\begin{abstract}
Rendering long text into images lets vision-language models (VLMs) read it with several times fewer tokens than the text itself.
This optical compression spends resolution uniformly.
Yet strings that must be copied exactly, such as identifiers and hashes, are misread into well-formed but wrong answers far more often than prose, and they fail more often as contexts grow.
We ask where the pixels should go.
The renderer, unlike the model, sees the text, so it can flag high-entropy strings and draw them larger while drawing prose smaller, within the vision-token budget of uniform rendering.
This allocation is training-free, query-independent, and needs only black-box access to the VLM.
Across five commercial and open-weight VLMs accessed through APIs, it raises RULER needle recall over uniform rendering by 0.55 and 0.43 at the two densest budgets (95\% CIs [0.48, 0.62] and [0.27, 0.56]).
At 32K--128K characters it raises recall by 0.48 while using up to 27\% fewer billed tokens.
Enlarging random words, or boxing the needles without enlarging them, does not help, so the gain comes from legibility of the targeted strings.
The gain shrinks as the budget relaxes and is no larger than a font search at 10\,px.
On LongBench, allocation doubles the recall of naturally occurring identifiers that the detector flags (0.23 to 0.48) without changing QA accuracy relative to uniform rendering; strings the detector misses gain nothing.
For flagged targets, verbatim anchors from prior work are as accurate on the allocated image as on a uniform one, and cheaper wherever allocation shrinks the image.
All image arms remain well below plain text on QA: allocation makes optical compression safer for identifiers, not a substitute for text.
\end{abstract}
